% from ega0-22-fr.tex
\oldpage[0\textsubscript{IV-1}]{188}

un $k$-homomorphisme (19.3.11), et l'on a donc déjà
$\operatorname{gr}^0(u)=v^0$ et $\operatorname{gr}^1(u)=v^1$. Or, on a le
diagramme commutatif
\[
\xymatrix@C=5em@R=3em{
  \mathbf S_K^*(\mathfrak m/\mathfrak m^2) \ar[r] \ar[d]_{w} &
  \operatorname{gr}_{\mathfrak m}^*(A) \ar[d]^{\operatorname{gr}^*(u)} \\
  \mathbf S_K^*(\mathfrak m'/\mathfrak m'^2) \ar[r] &
  \operatorname{gr}_{\mathfrak m'}^*(A')
}
\]
où les flèches horizontales sont les homomorphismes canoniques, et $w$
provient de $v^0$ et $v^1$~; comme $v^0$ et $v^1$ sont bijectifs, il en est
de même de $w$, et puisque $A$ et $A'$ sont des $k$-algèbres formellement
lisses, on déduit de (22.2.2, (i)) que $\operatorname{gr}(u)$ est bijectif~;
mais comme $A$ et $A'$ sont séparés et complets, cela entraîne que $u$
lui-même est \emph{bijectif} (Bourbaki, \emph{Alg. comm.}, chap.~III, §~2,
n\textsuperscript{o}~8, cor.~3 du th.~1).

(ii) Reste à montrer que l'image de (22.2.4.1) est l'ensemble des
homomorphismes \emph{injectifs} de $\Upsilon_{K/k}$ dans $V$~; on sait déjà
par (22.2.2, (ii)) qu'elle est contenue dans cet ensemble~; il suffira alors
de prouver le lemme suivant.
\end{proof}

\begin{lemma}[22.2.5.2]
\label{0.22.2.5.2-fr}
Pour tout $K$-homomorphisme $h:\Upsilon_{K/k}\to V$, il existe une
$k$-algèbre $A$ qui est un anneau local noethérien régulier et complet,
d'idéal maximal $\mathfrak m$, de corps résiduel $K$, et un $K$-isomorphisme
$\beta:\mathfrak m/\mathfrak m^2\xrightarrow{\sim}V$ tels que
$h=\beta\circ\chi_{A/k}$.
\end{lemma}

\begin{proof}
En effet, si ce lemme est prouvé, le fait que $A$ est régulier entraîne que
la condition (i) de (22.2.2) est satisfaite (17.1.1)~; si en outre $h$ est
injectif, (22.2.2) montrera que $A$ est une $k$-algèbre formellement lisse
dont la classe aura $h$ pour image (par contre si $h$ n'est pas injectif, la
$k$-algèbre $A$ dont (22.2.5.2) prouve l'existence n'est pas formellement
lisse).

Pour prouver (22.2.5.2), soient $B$ la $K$-algèbre $\mathbf S_K^*(V)$,
$\mathfrak n$ l'idéal d'augmentation de $B$ et $A=\widehat B$ la
$K$-algèbre complétée de $B$ pour la topologie $\mathfrak n$-préadique (de
sorte que $A$ est isomorphe à l'algèbre de séries formelles
$K[[T_1,\ldots,T_n]]$ si $n=\operatorname{rg}_KV$). En vertu de (20.6.11)
il existe sur $A/\mathfrak m^2$ (où $\mathfrak m=\mathfrak nA$) une structure
de $k$-extension de $K$ par $V$ (\emph{distincte} en général de la structure
$k$-triviale définie par l'injection canonique $k\to K\to A$) telle que $h$
soit l'homomorphisme caractéristique de cette extension. Si
$f:k\to A/\mathfrak m^2$ est l'homomorphisme structural, le fait que $k$
soit séparable sur le corps premier permet (19.6.1) de factoriser $f$ en
$k\xrightarrow{g}A\to A/\mathfrak m^2$, et la $k$-algèbre $A$ ainsi définie
répond à la question (20.6.22.2).
\end{proof}

\begin{theorem}[22.2.6]
\label{0.22.2.6-fr}
Soient $k$ un corps, $K$ une extension de $k$. Pour qu'il existe une
$k$-algèbre locale noethérienne $A$, d'idéal maximal $\mathfrak m$, telle que
$A/\mathfrak m=K$ et que $A$ soit formellement lisse (pour sa topologie
$\mathfrak m$-préadique), il faut et il suffit que $\Upsilon_{K/k}$ soit un
$K$-espace vectoriel de rang fini. La $k$-algèbre $A$ est alors déterminée
(à un $k$-isomorphisme près donnant par passage aux quotients l'identité
$K\to K$) par sa dimension, qui est soumise à la seule condition d'être
$\geqslant\operatorname{rg}_K\Upsilon_{K/k}$.
\end{theorem}

\begin{proof}
Comme $\mathfrak m/\mathfrak m^2$ est de rang fini sur $K$, la nécessité de
la condition résulte de (22.2.2, (ii))~; inversement, (22.2.5) montre que la
condition est suffisante et que l'on peut prendre pour
$\mathfrak m/\mathfrak m^2$ un $K$-espace vectoriel de rang quelconque
$\geqslant\operatorname{rg}_K(\Upsilon_{K/k})$.
\end{proof}

